% abstract (main.tex 9-15), introduction.tex lines 1-26 and 279-345, complex.tex lines 1-56 (the table)
\begin{abstract}
We disprove Sidorenko's conjecture with a bipartite graph on \(35\) vertices
and \(66\) edges: its homomorphism density in some finite simple graph is
smaller than the conjectured lower bound. The same connected graph also
disproves the forcing conjecture: at one fixed density, asymptotically
matching the edge and pattern densities does not imply quasirandomness.
\end{abstract}
\section{Introduction}\label{sec:introduction}

For finite simple undirected graphs \(H\) and \(G\) with
\(|V(G)|\ge1\), write
\[
t(H,G)=\frac{\Hom(H,G)}{|V(G)|^{|V(H)|}},
\qquad
p(G)=\frac{2|E(G)|}{|V(G)|^2},
\]
where \(\Hom(H,G)\) counts every edge-preserving map, including
noninjective maps. Sidorenko's conjecture asserts that
\begin{equation}\label{eq:sidorenko}
t(H,G)\ge p(G)^{|E(H)|}
\end{equation}
whenever \(H\) is bipartite and has at least one edge.
We disprove this conjecture.

\begin{theorem}\label{thm:main}
Let \(H\) be the incidence graph of the \(22\) triples on \(13\) points
in Table~\ref{tab:complex}. This bipartite graph has \(35\) vertices
and \(66\) edges. There is a finite simple undirected graph \(G\)
with at least one edge such that
\[
t(H,G)<p(G)^{66}.
\]
\end{theorem}
\begin{corollary}[Amplification and injective supersaturation]
\label{cor:amplification}
For the same graph \(H\) as in Theorem~\ref{thm:main}, the following
hold.
\begin{enumerate}[label=(\roman*)]
\item For every \(c>0\), there is a finite simple undirected graph
\(F\) with \(p(F)>0\) such that
\[
t(H,F)<c\,p(F)^{66}.
\]
\item There are \(p_*\in(0,1)\), \(\delta_*>0\), and finite simple
undirected graphs \(F_s\) of orders \(n_s\to\infty\) for which
\(d_s=|E(F_s)|/\binom{n_s}{2}\to p_*\) and, for all sufficiently
large \(s\),
\[
\operatorname{Inj}(H,F_s)
\le(1-\delta_*)(n_s)_{35}d_s^{66}.
\]
\end{enumerate}
\end{corollary}

The proof in Section~\ref{sec:consequences} uses tensor powers and
balanced blow-ups of a single host from Theorem~\ref{thm:main}.

Part (ii) contradicts the injective Erd\H{o}s--Simonovits
random-graph benchmark described in
\cite[Section~1]{conlonfoxsudakov2010}: the sequence has one positive
limiting density, so it eventually satisfies \(d_s>n_s^{-\gamma}\)
for every fixed \(\gamma>0\), while retaining a fixed deficit.

\subsection*{A consequence for the forcing conjecture}

Fix \(0<p<1\). A graph \(F\) with \(m\) edges is \emph{\(p\)-forcing}
if every sequence of finite simple graphs \(G_n\) with
\(|V(G_n)|\to\infty\), \(p(G_n)\to p\), and
\(t(F,G_n)\to p^m\) is quasirandom of density \(p\): for every fixed
finite simple graph \(J\), one has
\(t(J,G_n)\to p^{|E(J)|}\). It is \emph{forcing} if it is
\(p\)-forcing for every \(p\in(0,1)\). The classical edge-and-four-cycle
characterization of quasirandomness is due to Chung, Graham, and
Wilson~\cite{chunggrahamwilson1989}. Skokan and Thoma studied which
bipartite patterns can replace the four-cycle~\cite{skokanthoma2004}.
Using this asymptotic formulation and following their question,
Conlon, Fox, and Sudakov formulate the forcing conjecture: a graph is forcing if and
only if it is bipartite and contains a cycle
\cite[Section~1, equation~(2) and Conjecture~2]{conlonfoxsudakov2010}.

For a symmetric \([0,1]\)-valued kernel \(W\) on a probability space,
write
\[
t(F,W)=\E\prod_{uv\in E(F)}W(X_u,X_v),
\]
where the \(X_v\) are independent samples. In particular,
\(t(K_2,W)=\E W\). We identify each finite probability kernel below
with its step graphon on \([0,1]\).

\begin{corollary}[Failure of the forcing conjecture]
\label{cor:forcing}
For the same connected bipartite graph \(H\) of
Theorem~\ref{thm:main}, there are \(p\in(0,1)\) and a nonconstant
graphon \(W\) such that
\[
t(K_2,W)=p,\qquad t(H,W)=p^{66}.
\]
Consequently, \(H\) is not \(p\)-forcing at this density and hence is
not forcing, despite being bipartite and containing a cycle.
\end{corollary}
\section{The finite complex}
\label{sec:complex}

We specify the fixed triple system and certify the properties used in
both the kernel construction and the singular-tail bound.

Let $I=\{0,\ldots,12\}$ and $J=\{0,\ldots,21\}$. Associate to
each $j\in J$ the three-element subset of $I$ in Table~\ref{tab:complex};
we call these subsets \emph{faces} and also denote them by $j$.
Let $E$ be the set of unordered point pairs contained in a face.
The graph used throughout the paper is the incidence graph
\[
 H=(I\sqcup J,\mathscr C),\qquad
 \mathscr C=\{(i,j):i\in j\}.
\]
The copies of $I$ and $J$ in this disjoint union are distinct even when
their numerical labels agree. We call an incidence $(i,j)$ a
\emph{corner}. Thus $H$ has $35$ vertices and $66$ edges.

\begin{table}[htbp]
\centering
\small
\setlength{\tabcolsep}{7pt}
\begin{tabular}{c@{\qquad}ccc@{\qquad}ccc@{\qquad}c}
\toprule
Face&\multicolumn{3}{c}{Points}&\multicolumn{3}{c}{Opposite face}&Bits\\
 $j$&$i_1$&$i_2$&$i_3$&$j_{12}$&$j_{13}$&$j_{23}$&\\
\midrule
0 &0&1&3   &1&2&4   &000\\
1 &0&1&9   &0&3&7   &010\\
2 &0&2&3   &3&0&8   &001\\
3 &0&2&9   &2&1&9   &110\\
4 &1&3&10  &0&5&11  &001\\
5 &1&7&10  &6&4&16  &000\\
6 &1&7&12  &5&7&18  &100\\
7 &1&9&12  &1&6&21  &110\\
8 &2&3&4   &2&9&10  &101\\
9 &2&4&9   &8&3&12  &100\\
10&3&4&11  &8&11&13 &111\\
11&3&10&11 &4&10&17 &101\\
12&4&8&9   &13&9&21 &101\\
13&4&8&11  &12&10&19&011\\
14&5&6&7   &15&16&18&000\\
15&5&6&11  &14&17&19&110\\
16&5&7&10  &14&17&5 &010\\
17&5&10&11 &16&15&11&111\\
18&6&7&12  &14&20&6 &001\\
19&6&8&11  &20&15&13&010\\
20&6&8&12  &19&18&21&011\\
21&8&9&12  &12&20&7 &111\\
\bottomrule
\end{tabular}
\caption{The faces of the complex. The entry $j_{ab}$ is the other
face containing $\{i_a,i_b\}$. Bit positions are read from left to right.}
\label{tab:complex}
\end{table}
